\documentclass[11pt,a4paper]{article}
\usepackage[T1]{fontenc}
\usepackage[utf8]{inputenc}
\usepackage[a4paper,margin=2.5cm]{geometry}
\usepackage[UKenglish]{babel}
\usepackage{hyperref}
\hypersetup{hidelinks}
\setlength{\parskip}{0.6em}
\setlength{\parindent}{0pt}

\begin{document}

\noindent 30 April 2026

\bigskip

\noindent Editor-in-Chief\\
Proceedings of the Institution of Mechanical Engineers,\\
Part D: Journal of Automobile Engineering\\
SAGE Publishing

\bigskip

\noindent Dear Editor,

We submit for your consideration the manuscript entitled ``BowTie Risk
Analysis of the Valve Train in a High-Performance VW AP~1.8 Cylinder
Head Operating at 10{,}000~rpm'' as an Original Article.

The manuscript proposes a quantitative-coupled BowTie procedure for
design-stage reliability assessment of the valve train of a
high-performance spark-ignition cylinder head, and applies it to a
top event of catastrophic significance for high-rpm engines: the loss
of follower contact between valve and cam at 10{,}000~rpm. Beyond the
qualitative bow diagram, each preventive barrier is associated with a
deterministic indicator computed from the analytical sizing of the
intake valve and the dual chrome--silicon valve springs (dynamic
follower-contact ratio, natural-frequency margin against cam
harmonics, coil-bind clearance, and Goodman fatigue factor of safety),
with literature-anchored pass criteria. The procedure is qualitatively
validated against four documented valve-train and engine-component
failure cases from the recent literature, and a concrete experimental
validation roadmap --- spring-rig modal verification, motored-engine
speed sweep, and inner-spring fatigue campaign, with explicit
acceptance criteria --- is specified for empirical confirmation. The
case study is the Volkswagen AP~1.8 platform --- a base engine with a
wide installed base in the Brazilian automotive market --- for which a
high-performance cylinder head is currently being developed within a
Master's research project at the Universidade Cat\'olica de Petr\'opolis.

We believe \emph{Proceedings of the IMechE, Part D: Journal of
Automobile Engineering} is the most appropriate venue because (i)~the
journal has a long publication history of design and reliability
analyses of internal-combustion-engine subsystems, including
valve-train and cam-train dynamics; (ii)~the proposed
quantitative-coupling layer addresses a recurring gap between
qualitative risk diagrams and engineering design verification, which
is of direct interest to the journal's automotive-engineering
readership; and (iii)~the application domain --- valve-train dynamics
at extreme rotational speed in a passenger-car-derived platform ---
aligns with the journal's traditional scope.

The work is co-authored by Darlan Costa Porto (corresponding author,
Master's student in the Engineering Systems programme), Heitor
Oliveira Gon\c{c}alves (faculty of the same programme, ORCID
0000-0002-5866-2213), Florian Pradelle (PUC-Rio collaborator), and
Jos\'e Cristiano Pereira (advisor of the underlying dissertation).
Author contributions are detailed in the CRediT statement at the end
of the manuscript.

The accompanying manuscript PDF was prepared in a generic LaTeX
double-column template. The LaTeX source can be ported to the SAGE
journal style file (\texttt{sagej.cls}) on acceptance, in line with
the journal's usual production workflow. The manuscript is original,
has not been published elsewhere, and is not under consideration by
any other journal. All authors have approved the submission, and the
authors declare no competing financial interests.

We respectfully suggest the following independent reviewers as
potentially appropriate. None has co-authored work cited in this
manuscript with the present authors, and none is affiliated with the
Universidade Cat\'olica de Petr\'opolis. Institutional contact details
can be located on each researcher's institutional page.

\begin{enumerate}
\item Prof.\ Chiara Soffritti --- Department of Engineering, University
of Ferrara, Italy. Co-author of a 2018 paper on valve-train wear in
four-cylinder diesel engines (DOI 10.1016/j.engfailanal.2018.06.022),
directly comparable to the present case study.

\item Dr.\ Jos\'e Antonio V\'elez Godi\~no --- Departamento de
Ingenier\'ia Energ\'etica, Universidad de Sevilla, Spain. First author
of a 2019 paper on overhead valve-train failure in urban buses
(DOI 10.1016/j.engfailanal.2018.11.003); competent in field-data
interpretation of valve-train degradation.

\item Prof.\ Faisal Khan --- Mary Kay O'Connor Process Safety Center,
Texas A\&M University, USA. Co-author of the BowTie--Bayesian-network
mapping (Khakzad, Khan and Amyotte, 2013,
DOI 10.1016/j.psep.2012.01.005); senior authority on barrier-based
risk methods.

\item Prof.\ Frank Guldenmund --- Section Safety and Security Science,
TU Delft, Netherlands. Co-author of the canonical BowTie review
(de Ruijter and Guldenmund, 2016, DOI 10.1016/j.ssci.2016.03.001).

\item Dr.\ Jonas Aust --- Department of Mechanical Engineering,
University of Canterbury, New Zealand. First author of the BowTie
methodology paper for aircraft-engine borescope inspection
(\emph{Aerospace}, 2019, DOI 10.3390/aerospace6100110); experienced
in BowTie applied to engine maintenance.
\end{enumerate}

The final selection of reviewers is, of course, at the editor's
discretion.

We look forward to your editorial decision and remain available for
any clarification.

\bigskip

\noindent Yours sincerely,

\bigskip\bigskip

\noindent Darlan Costa Porto\\
Corresponding author\\
Programa de Mestrado em Sistemas de Engenharia\\
Universidade Cat\'olica de Petr\'opolis\\
Petr\'opolis, RJ, Brazil\\
\href{mailto:darlan.42540004@ucp.br}{darlan.42540004@ucp.br}

\end{document}
